\section*{个人简介}
刘天成是一名计算媒体研究者与艺术实践者，研究方向横跨创意人工智能、多模态交互与文化语境下的计算方法。他关注机器如何理解、感知并参与人类创作实践，尤其聚焦中国书法等具身文化形式。
他的博士研究围绕创作实践中的三个相互关联层面展开。第一层面关注机器如何分析和解释完成作品中的审美结构，包括章法、形态、平衡与风格偏好。第二层面关注机器如何感知创作过程中常常不可见的过程性信息，例如手势、压力、节奏和其他隐性的微动作。第三层面关注 AI 系统如何作为协作者参与创作实践，支持共同创作、表演与表达性交互。
在此基础上，他希望进一步探索计算方法作为理解、表达与协作媒介的可能性，拓展人机如何在文化与艺术语境中共同感知、创造与想象。
